\documentclass[aspectratio=169,11pt]{beamer}

\usetheme{Madrid}
\usecolortheme{default}
\usefonttheme{professionalfonts}
\setbeamertemplate{navigation symbols}{}
\setbeamertemplate{footline}[frame number]

\usepackage{amsmath, amssymb, booktabs, array}

\title{Incentives, Contracts, Reciprocity, and Workplace Behavior}
\subtitle{Behavioral Labor -- Week 5}
\author{Teaching deck draft}
\date{}

\begin{document}

\begin{frame}
  \titlepage
\end{frame}

\begin{frame}{Central question}
\begin{itemize}
    \item How do behavioral mechanisms reshape incentives and effort inside the employment relationship?
    \item Week 5 moves inside firms: pay, contracts, monitoring, supervision, and evaluation.
    \item The core issue is not whether incentives matter.
    \item The issue is how workers interpret incentives and which margin actually moves.
\end{itemize}
\end{frame}

\begin{frame}{Bridge from Weeks 2--4}
\small
\begin{tabular}{p{0.26\linewidth} p{0.60\linewidth}}
\toprule
Earlier object & Week 5 workplace shift \\
\midrule
Nonstandard preferences & Reciprocity, reference points, fairness, and loss framing at work \\
Beliefs and expectations & Workers infer what the firm values and what supervisors reward \\
Attention and complexity & Contracts, bonuses, and monitoring dashboards must be understood \\
Policy and firm design & Employers choose pay rules, messages, visibility, and evaluation systems \\
\bottomrule
\end{tabular}
\end{frame}

\begin{frame}{Workplace behavior as behavioral labor}
\small
\begin{tabular}{p{0.25\linewidth} p{0.31\linewidth} p{0.28\linewidth}}
\toprule
Contract object & Behavioral channel & Labor outcome \\
\midrule
Gift or wage gesture & reciprocity and fairness & extra work or productivity \\
Bonus frame & reference point and loss aversion & target response or gaming \\
Monitoring regime & visibility and salience & compliance or measured effort \\
Supervisor rating & subjective evaluation & influence activity or output \\
Performance pay & heterogeneous response & output, sorting, earnings dispersion \\
\bottomrule
\end{tabular}

\vspace{0.6em}
The empirical question: what changed, which margin moved, and what remains hidden?
\end{frame}

\begin{frame}{Benchmark effort and contract problem}
\[
\max_{e \ge 0} \; u\!\left(w(y(e))\right) - c(e)
\]

\small
\begin{tabular}{p{0.18\linewidth} p{0.68\linewidth}}
\toprule
Object & Interpretation \\
\midrule
$e$ & workplace effort \\
$y(e)$ & measured output or task completion \\
$w(y(e))$ & compensation schedule tied to output \\
$c(e)$ & effort cost \\
Benchmark & firms change pay rules and workers move effort \\
\bottomrule
\end{tabular}
\end{frame}

\begin{frame}{Behavioral contracts plus monitoring}
\[
\max_{e \ge 0} \; u\!\left(w(y(e),m)\right) - c(e)
 + B\!\left(g_i, r_i, f_i, s_i, m_i; e\right)
\]

\begin{itemize}
    \item $m$: monitoring regime or information technology.
    \item $g_i$: perceived generosity or employer effort.
    \item $r_i$: reference point for pay, bonus, target, or treatment.
    \item $f_i$: framing and contract presentation.
    \item $s_i$: social preferences, reciprocity, fairness, and workplace norms.
\end{itemize}
\end{frame}

\begin{frame}{Incentives plus monitoring}
\small
\begin{tabular}{p{0.24\linewidth} p{0.31\linewidth} p{0.29\linewidth}}
\toprule
Monitoring role & Possible benefit & Possible cost \\
\midrule
Make output visible & incentives bite on productive effort & workers optimize the metric \\
Increase feedback & learning and task correction & stress or crowd-out \\
Signal managerial attention & relational treatment or accountability & surveillance and mistrust \\
Support evaluation & better information for pay or promotion & influence activity and bias \\
\bottomrule
\end{tabular}

\vspace{0.7em}
Monitoring is part of incentive design, not a neutral add-on.
\end{frame}

\begin{frame}{Reciprocity and gift exchange}
\begin{itemize}
    \item Core object: workers may reciprocate employer generosity.
    \item Kube-Marechal-Puppe: gifts and wage gestures in workplace effort.
    \item DellaVigna-List-Malmendier-Rao: social preferences and gift exchange at work.
    \item Key diagnostic: does the treatment raise extra work, productivity, or both?
    \item Do not stop at ``gift exchange exists.'' Ask what margin moved.
\end{itemize}
\end{frame}

\begin{frame}{Reference dependence and loss framing}
\[
U(e) = u\!\left(w(y(e))\right) - c(e)
  + \lambda \min\{0,\, b(e)-r\}
\]

\begin{itemize}
    \item A bonus can be framed as a gain or as something at risk of being lost.
    \item Loss framing can increase urgency around a target.
    \item But target response is not the same as total performance.
    \item Frontier question: when does loss framing motivate, and when does it induce gaming?
\end{itemize}
\end{frame}

\begin{frame}{Subjective evaluation and supervision}
\[
\max_{e,a \ge 0} \; u\!\left(w\big(z(e), \hat q^S(e,a)\big)\right)
  - c(e) - \psi(a)
\]

\small
\begin{tabular}{p{0.16\linewidth} p{0.70\linewidth}}
\toprule
Object & Interpretation \\
\midrule
$z(e)$ & objective productive output \\
$\hat q^S(e,a)$ & supervisor's subjective assessment \\
$a$ & influence activity or supervisor-facing effort \\
Risk & workers please evaluators instead of improving production \\
\bottomrule
\end{tabular}
\end{frame}

\begin{frame}{Performance pay, heterogeneity, and sorting}
\[
k_i^\star \in \arg\max_{k \in \mathcal{K}}
 \; \mathbb{E}\!\left[U_i\!\left(w_k(y_i(e),m), e\right)\right]
\]

\begin{itemize}
    \item Performance-pay estimates can mix treatment and selection.
    \item Workers differ in monetary responsiveness, risk tolerance, reciprocity, and monitoring aversion.
    \item Contract choice can change the workforce, not only effort of incumbents.
    \item Average effects can hide large gains, crowd-out, or unequal responses by type.
\end{itemize}
\end{frame}

\begin{frame}{Productivity versus extra work versus gaming}
\small
\begin{tabular}{p{0.26\linewidth} p{0.28\linewidth} p{0.30\linewidth}}
\toprule
Margin & Looks like success when... & Risk \\
\midrule
Productive effort & valuable output rises & hard to measure fully \\
Extra work & time or tasks increase & activity is not value \\
Gaming & target attainment rises & metric improves, production does not \\
Influence activity & ratings improve & supervisor-facing effort crowds out work \\
Sorting & high responders enter & treatment effect is overstated \\
\bottomrule
\end{tabular}
\end{frame}

\begin{frame}{Frontier: digital and algorithmic monitoring}
\begin{itemize}
    \item Remote work and platform work bundle pay, monitoring, feedback, and discipline.
    \item Digital tools can reveal output, process, time, location, communication, and compliance.
    \item Behavioral questions: visibility, trust, stress, gaming, reciprocity, and privacy costs.
    \item Identification question: does monitoring change productivity, measured effort, evaluation, or worker sorting?
\end{itemize}
\end{frame}

\begin{frame}{Methods map}
\scriptsize
\begin{tabular}{p{0.28\linewidth} p{0.31\linewidth} p{0.27\linewidth}}
\toprule
Design & Identifying object & Main caution \\
\midrule
Gift experiment & reciprocity and generosity response & kindness versus surprise \\
Framing experiment & reference dependence and loss response & target attainment versus value \\
Monitoring experiment & observation and supervision effects & visibility versus motivation \\
Subjective evaluation & influence activity and ratings & true quality versus relationship capital \\
Real-effort contracts & relative motivator strength & workplace external validity \\
Contract choice & treatment versus sorting & endogenous selection \\
\bottomrule
\end{tabular}
\end{frame}

\begin{frame}{Week 5 lab}
\begin{itemize}
    \item Reproduce: synthetic gift-exchange workplace factbook.
    \item Diagnose: behavioral object, contract element, labor margin, and identification claim.
    \item Transfer: monitoring intensity and productivity versus compliance.
    \item Frontier extension: subjective evaluation and influence activity.
\end{itemize}
\end{frame}

\begin{frame}{Bridge to Week 6}
\begin{itemize}
    \item Week 5: firms govern behavior inside the employment relationship.
    \item Pay, monitoring, gifts, loss frames, and evaluation change effort composition.
    \item Week 6 widens the lens to identity, norms, fairness, and labor-market allocation.
    \item The bridge is social meaning: contracts are interpreted inside organizations and markets.
\end{itemize}
\end{frame}

\begin{frame}{Takeaways}
\begin{itemize}
    \item Incentives work through pay, interpretation, monitoring, and evaluation.
    \item Extra work, productivity, gaming, and influence activity are different outcomes.
    \item Stronger incentives are not automatically better workplace design.
    \item Credible evidence names the behavioral object, labor margin, outcome metric, and identification claim.
\end{itemize}
\end{frame}

\end{document}
